\documentclass[11pt,a4paper]{article}
\usepackage[margin=1in]{geometry}
\usepackage{amsmath}
\usepackage{hyperref}
\usepackage{enumitem}
\usepackage{graphicx}
\usepackage{xcolor}
\usepackage{longtable}
\usepackage{array}

% -----------------------------------------------------
% bvraghav-tiet-ucs503p-article.sty
% -----------------------------------------------------

% Variable: Lab Instructor
\makeatletter \newcommand{\BvrSetLabInstructor}[1]{%
  \gdef\Bvr@LabInstructor{#1}}
% default
\newcommand{\Bvr@LabInstructor}{Dr. Paramveer %
  Sidhu}
% public accessor
\newcommand{\BvrLabInstructorValue}{\Bvr@LabInstructor}
\makeatother

% Add subtitle
\usepackage{titling}
% -- Set title bold
\pretitle{\begin{center}\bfseries\fontsize{18bp}{18bp}\selectfont}
\makeatletter
\newcommand{\BvrSubtitle}[1]{%
  \posttitle{%
    \par\end{center}
    \begin{center}\large#1\end{center}
    \vskip 0.5em}%
}
\makeatother

% Hints in the section title(s)
\newcommand{\BvrHint}[1]{%
  {\normalfont\normalsize\itshape\hfill #1}}

% Requirement-ID macro: \ReqID{FR-1}
\newcommand{\ReqID}[1]{\texttt{#1}}

% -----------------------------------------------------

\title{Aegis --- An Encrypted ``Dead Man's Switch'' Legacy Vault}

\BvrSetLabInstructor{Dr. Paramveer Sidhu}

\BvrSubtitle{%
  Software Requirements Specification (SRS) for UCS503P \\
  Submitted to: \BvrLabInstructorValue \\ \bfseries
  Thapar Institute of Engineering and Technology%
}

\author{
  Pratham Arora, Krishna Pandey, Nipun Behl (COE)
}

\date{\today}

\begin{document}
\maketitle

\begin{abstract}
  \noindent Aegis is an encrypted ``dead man's switch'' legacy vault. An
  \emph{Owner} deposits encrypted files or messages into a \emph{vault} and
  designates $N$ \emph{trustees} with a reconstruction \emph{threshold} $K$.
  The system periodically requires the Owner to \emph{check in}; if a check-in
  is missed and a configured \emph{grace period} lapses, the decryption key
  --- pre-split using Shamir's Secret Sharing so that any $K$ of $N$ shares
  reconstruct it and any $K-1$ reveal nothing --- is distributed to the
  trustees, who together reconstruct the key and open the vault. This document
  specifies the software requirements for version~1 (v1) of Aegis.
\end{abstract}

\section{Problem Statement \BvrHint{\ldots the gap this project addresses}}
Individuals routinely hold digital secrets that they cannot safely disclose while
alive, yet which must not be lost when they die or become permanently
incapacitated: banking and brokerage credentials, cryptocurrency private keys,
password-manager master passwords, encrypted archives, business continuity
information, and personal messages intended for specific recipients. The value of
such data is entirely conditional on \emph{access}, and access is precisely what
is destroyed by the owner's absence.

Existing practice offers two options, and both are unsatisfactory:
\begin{enumerate}[leftmargin=0.9cm,label=(\alph*)]
  \item \textbf{Disclose now.} The owner shares the secret with a trusted person
    or writes it into a document held by a custodian. This converts a
    confidentiality problem into a trust problem that must hold indefinitely: the
    secret is readable from the moment it is shared, and its safety depends on the
    continued honesty, competence, and availability of a single party.
  \item \textbf{Disclose never.} The owner tells no one. Confidentiality is
    preserved perfectly, and the data is irrecoverably lost on the owner's death.
\end{enumerate}

Conventional instruments such as sealed wills and escrow arrangements sit within
option~(a) and inherit its central weakness --- a \textbf{single point of
failure}. One document, one custodian, and one occasion on which it is opened: if
the document is destroyed, the custodian is compromised or unreachable, or the
release is contested, the owner's intent is not carried out. Furthermore, these
instruments are triggered by an external legal or administrative process rather
than by the absence of the owner, and are therefore unsuitable for digital assets
that require timely release.

The problem this project addresses is therefore:

\begin{quote}
  \itshape How can a system release a secret to designated recipients when, and
  only when, its owner has become unresponsive --- without any single party,
  including the system operator, being able to read that secret beforehand, and
  without the release depending on any single custodian?
\end{quote}

A solution must satisfy three properties simultaneously, and it is their
combination that makes the problem non-trivial:
\begin{itemize}[leftmargin=0.7cm]
  \item \textbf{Conditional release.} Disclosure must be triggered automatically
    by the owner's sustained non-response, not by a human decision.
  \item \textbf{Distributed trust.} No single recipient may hold enough material
    to open the vault alone, so that neither the loss nor the betrayal of one
    party is sufficient to defeat the system.
  \item \textbf{Temporal safety.} The system must never disclose early. A
    premature release is unrecoverable --- the secret cannot be un-revealed ---
    and is therefore treated in this specification as the catastrophic failure
    mode (\ReqID{NFR-REL-1}).
\end{itemize}

\textbf{Aegis} addresses this problem by combining a \emph{dead man's switch} ---
a periodic check-in whose absence, rather than whose presence, is the trigger ---
with \emph{Shamir's Secret Sharing}, which splits the vault's decryption key into
$N$ shares such that any $K$ reconstruct it and any $K-1$ reveal nothing about it
in an information-theoretic sense. The former provides conditional release; the
latter eliminates the single custodian.

\section{Purpose \BvrHint{\ldots why this document exists}}
This Software Requirements Specification (SRS) defines the functional and
non-functional requirements of \textbf{Aegis}, an encrypted legacy vault built
around a time-triggered release mechanism. It is the authoritative reference
for the design, implementation, and evaluation of v1 across the UCS503P
project timeline. The intended readers are the lab instructor/evaluator (as the
acceptance authority) and the project team (as the developers). Each functional
requirement carries a traceable identifier (\ReqID{FR-\textit{n}}) and each
non-functional requirement a category identifier (\ReqID{NFR-\textit{cat}-\textit{n}})
so that later design, code, and tests can cite the requirement they satisfy.

\section{Scope \BvrHint{\ldots what v1 is, and is not}}
\subsection{In scope (v1)}
Aegis v1 is a web application with a Python/FastAPI backend and a React
frontend. It provides:
\begin{itemize}[leftmargin=0.7cm]
  \item account registration and authentication for Owners;
  \item creation of a single vault per Owner and upload of an encrypted
    payload;
  \item configuration of a check-in interval and a grace period;
  \item designation of $N$ trustees and a threshold $K$ ($1 \le K \le N$);
  \item a scheduler that issues check-in prompts and, on missed check-in past
    the grace period, distributes key shares to trustees;
  \item a one-action check-in confirmation for the Owner;
  \item trustee-side submission of $K$ shares to reconstruct the key and
    decrypt the payload.
\end{itemize}
The algorithmic core is a hand-written implementation of \emph{Shamir's Secret
Sharing} (SSS) over a finite field, together with symmetric encryption of the
payload under the shared key.

\subsection{Non-goals \BvrHint{\ldots explicit exclusions for v1}}
The following are \textbf{explicitly out of scope} for v1 and must not be
assumed by design or evaluation:
\begin{enumerate}[leftmargin=0.9cm,label=\textbf{NG-\arabic*.}]
  \item \textbf{No SMS.} All notifications are delivered by email (SMTP) only;
    no SMS or telephony channel is provided.
  \item \textbf{No mobile app.} Aegis v1 is a responsive web application only;
    no native Android/iOS client is built.
  \item \textbf{Single vault per user.} Each Owner has at most one vault in v1;
    multi-vault management is deferred.
\end{enumerate}

\section{Glossary \BvrHint{\ldots shared vocabulary}}
\begin{longtable}{>{\raggedright\arraybackslash}p{0.24\textwidth}%
    >{\raggedright\arraybackslash}p{0.68\textwidth}}
  \textbf{Term} & \textbf{Definition} \\ \hline \endhead
  \textbf{Vault} & The encrypted container owned by a single Owner, holding one
    payload plus its release configuration and lifecycle state. \\
  \textbf{Payload} & The Owner's secret data (files and/or messages), stored
    only as ciphertext. \\
  \textbf{Share} & One piece of the split decryption key produced by Shamir's
    Secret Sharing; held by a single trustee. \\
  \textbf{Threshold ($K$)} & The minimum number of shares required to
    reconstruct the key; any $K$ of $N$ suffice, any $K-1$ reveal nothing. \\
  \textbf{Grace period} & The additional time granted after a missed check-in
    before the vault is released. \\
  \textbf{Check-in} & An explicit, confirmed action by which the Owner signals
    they are still present, resetting the release timer. \\
  \textbf{Owner} & The registered user who creates a vault, uploads the
    payload, and is prompted to check in. \\
  \textbf{Trustee} & A person designated by the Owner to receive a key share
    and participate in reconstruction after release. \\
  \textbf{Scheduler} & The time-triggered system actor that tracks deadlines
    and drives the vault through its lifecycle. \\
\end{longtable}

\section{System actors and models
  \BvrHint{\ldots see \texttt{docs/} diagrams}}
Aegis recognises three actors: the \textbf{Owner} and \textbf{Trustee} (human)
and the \textbf{Scheduler} (a time-triggered system actor that initiates
prompts and release without human action). The following models accompany this
SRS and are maintained as Mermaid diagrams in the project documentation:
\begin{itemize}[leftmargin=0.7cm]
  \item \textbf{Use-case diagram} --- \texttt{docs/diagrams/use-case.md}
    (actors Owner, Trustee, Scheduler).
  \item \textbf{Vault lifecycle state diagram} ---
    \texttt{docs/diagrams/vault-lifecycle.md}
    (Active $\rightarrow$ Warning $\rightarrow$ Grace $\rightarrow$ Released).
  \item \textbf{High-level architecture} ---
    \texttt{docs/diagrams/architecture.md}.
\end{itemize}

\section{Functional Requirements
  \BvrHint{\ldots traceable \ReqID{FR-\textit{n}}}}
\begin{description}[leftmargin=0cm,style=nextline]
  \item[\ReqID{FR-1} Registration \& authentication.] The system shall allow a
    visitor to register as an Owner with an email and password, and shall
    authenticate returning Owners. Credentials shall be verified before any
    vault operation. Passwords are never stored in plaintext (see
    \ReqID{NFR-SEC-3}).
  \item[\ReqID{FR-2} Vault creation \& payload upload.] An authenticated Owner
    shall be able to create their (single) vault and upload a payload (files
    and/or a text message). The payload shall be encrypted client- or
    server-side and persisted only as ciphertext (see \ReqID{NFR-SEC-1}).
  \item[\ReqID{FR-3} Timing configuration.] The Owner shall configure the
    \emph{check-in interval} (how often they must check in) and the
    \emph{grace period} (additional time after a missed check-in before
    release). Both are positive durations; the grace period may be zero only
    if explicitly chosen.
  \item[\ReqID{FR-4} Trustee designation \& threshold.] The Owner shall
    designate $N \ge 1$ trustees (by email) and set a threshold $K$ with
    $1 \le K \le N$. On confirmation, the system shall split the vault key into
    $N$ Shamir shares with threshold $K$.
  \item[\ReqID{FR-5} Scheduled check-in prompts.] The Scheduler shall, at each
    interval boundary, transition the vault to \emph{Warning} and prompt the
    Owner to check in via email (see \ReqID{NFR-PERF-2}).
  \item[\ReqID{FR-6} One-action check-in confirmation.] The Owner shall be able
    to confirm a check-in with a single action (one click on a tokenised link
    or button), which resets the release timer and returns the vault to
    \emph{Active}.
  \item[\ReqID{FR-7} Share distribution on expiry.] If no valid check-in is
    confirmed before the interval deadline \emph{and} the grace period have
    both elapsed, the Scheduler shall transition the vault to \emph{Released}
    and distribute one share to each trustee by email --- and shall do so
    \textbf{only then} (see \ReqID{NFR-REL-1}).
  \item[\ReqID{FR-8} $K$-of-$N$ reconstruction \& decryption.] After release,
    the system shall accept shares submitted by trustees and, upon receiving
    any $K$ valid shares, reconstruct the key and decrypt the payload for
    retrieval. Fewer than $K$ shares shall not permit decryption.
\end{description}

\section{Non-Functional Requirements
  \BvrHint{\ldots measurable \ReqID{NFR-\textit{cat}-\textit{n}}}}
\subsection{Reliability \BvrHint{\ldots the catastrophic failure mode}}
\begin{description}[leftmargin=0cm,style=nextline]
  \item[\ReqID{NFR-REL-1} No early release.] The system shall \textbf{never}
    release a vault (transition to \emph{Released} or distribute any share)
    before the check-in deadline \emph{and} the full grace period have elapsed
    with no confirmed check-in. \emph{Measure:} across an automated reliability
    suite of $\ge 1000$ simulated schedules --- including process restarts,
    timer resets, retries, and system-clock adjustments --- the count of
    releases occurring earlier than \texttt{deadline + grace} shall be exactly
    \textbf{zero}. This is the project's single most important requirement.
  \item[\ReqID{NFR-REL-2} Durable deadlines.] Pending release deadlines shall
    survive a process/host restart. \emph{Measure:} after a forced restart at
    an arbitrary point, every previously scheduled deadline is re-armed and
    fires no later than 60~s after its due time, with no deadline lost and none
    fired early.
  \item[\ReqID{NFR-REL-3} Exactly-once release.] Each vault shall be released
    at most once and each trustee shall receive at most one share.
    \emph{Measure:} no duplicate share emails under retry/concurrent execution
    in the reliability suite.
\end{description}

\subsection{Security \BvrHint{\ldots payload confidentiality}}
\begin{description}[leftmargin=0cm,style=nextline]
  \item[\ReqID{NFR-SEC-1} Payload never in plaintext.] The payload shall never
    be persisted in plaintext; only ciphertext plus non-secret metadata
    (filename, size, MIME type) may be stored. \emph{Measure:} an automated
    inspection of the database and file store finds no plaintext payload for
    any vault; the plaintext key exists only transiently in memory during
    encryption and reconstruction.
  \item[\ReqID{NFR-SEC-2} Threshold secrecy.] Any set of at most $K-1$ shares
    shall reveal no information about the key. \emph{Measure:} property tests
    confirm that all $\binom{N}{K-1}$ share subsets are statistically
    consistent with every possible secret (an intrinsic guarantee of SSS).
  \item[\ReqID{NFR-SEC-3} Credential protection.] Passwords shall be stored
    only as salted hashes using a memory-hard function (e.g. Argon2/bcrypt).
    \emph{Measure:} no reversible or plaintext credential appears in storage.
  \item[\ReqID{NFR-SEC-4} Transport security.] All client--server traffic shall
    be over HTTPS/TLS in any deployed environment.
\end{description}

\subsection{Performance, usability, and maintainability}
\begin{description}[leftmargin=0cm,style=nextline]
  \item[\ReqID{NFR-PERF-1} Responsiveness.] A check-in confirmation shall be
    processed within 500~ms at the 95th percentile under nominal load.
  \item[\ReqID{NFR-PERF-2} Prompt timeliness.] A due check-in prompt shall be
    dispatched within 60~s of its scheduled time.
  \item[\ReqID{NFR-USE-1} One-action check-in.] Confirming a check-in shall
    require exactly one deliberate user action and no re-authentication beyond
    the tokenised link (supports \ReqID{FR-6}).
  \item[\ReqID{NFR-MAINT-1} Code quality.] The codebase shall pass \texttt{ruff}
    with no errors and maintain unit-test coverage $\ge 80\%$ on the
    \texttt{crypto} and \texttt{scheduler} packages.
  \item[\ReqID{NFR-PORT-1} Storage portability.] The system shall run against
    SQLite (development) and PostgreSQL (later) without code changes, via the
    SQLAlchemy abstraction.
\end{description}

\section{Assumptions and Constraints}
\subsection{Assumptions}
\begin{itemize}[leftmargin=0.7cm]
  \item The Owner and every trustee have a working, monitored email address.
  \item Trustees are cooperative and, after release, are willing to combine
    their shares to open the vault.
  \item The server's scheduling clock is the authoritative time source; the
    Owner's timezone is recorded for display only.
  \item Email delivery (SMTP) is reasonably reliable; transient failures are
    retried.
\end{itemize}
\subsection{Constraints}
\begin{itemize}[leftmargin=0.7cm]
  \item \textbf{Platform:} Python~3 + FastAPI backend, SQLAlchemy persistence,
    APScheduler scheduling, SMTP email, React~+~Tailwind frontend.
  \item \textbf{Algorithmic:} Shamir's Secret Sharing is hand-written (finite-field
    polynomial evaluation + Lagrange interpolation), not taken from a
    third-party secret-sharing library; no HSM is used in v1.
  \item \textbf{Product:} the v1 non-goals (NG-1\ldots NG-3) hold --- email
    only, web only, one vault per Owner.
  \item \textbf{Process:} the system is developed by a three-member team over a 12-week UCS503P
    timeline with weekly, CI-backed increments.
\end{itemize}

\section{Requirement traceability \BvrHint{\ldots forward references}}
The functional and non-functional identifiers defined here are cited directly
in the module documentation (\texttt{code/*/README.md}) and will be cited by
the corresponding tests as they are written, giving a requirement
$\rightarrow$ module $\rightarrow$ test trace. The vault lifecycle referenced
by \ReqID{FR-5}--\ReqID{FR-7} and \ReqID{NFR-REL-1} is specified by the state
diagram in \texttt{docs/diagrams/vault-lifecycle.md}.

\end{document}
